\documentclass[11pt,a4paper]{article}
\usepackage[margin=25mm]{geometry}
\usepackage{amsmath,amssymb,amsthm,booktabs}
\newtheorem{proposition}{Proposition}
\title{Two O'Nan map counts from maximal generation}
\author{Human Compression project}\date{}
\begin{document}\maketitle
Let $G=O'N$, and retain the ordered-pair convention
\[
 X_L(p,q)=\{(r,s)\in L^2:|r|=p,\ |s|=q,\ |rs|=2\},
 \qquad x_L(p,q)=|X_L(p,q)|.
\]
We use the previously acquired exact character evaluations, the
complete maximal-subgroup classification of $G$, and the fixed
581-class subgroup roster. These are explicit borrowed premises.
The following criterion supplies the missing generation step for two
entries of Connor--Leemans, Table~1.

\begin{proposition}
Let $G$ be a finite nonabelian simple group. Suppose that, for each
maximal subgroup $M$ with $x_M(p,q)>0$, none of the maximal subgroups
of $M$ has order divisible by $\operatorname{lcm}(p,q)$. Then every
nongenerating pair in $X_G(p,q)$ generates a maximal subgroup of $G$
and has exactly one maximal overgroup.
\end{proposition}
\begin{proof}
Choose a maximal subgroup $M$ containing a nongenerating pair and
put $H=\langle r,s\rangle$. If $H$ were proper in $M$, it would lie
in a maximal subgroup of $M$. Its order would then be divisible by
both $p$ and $q$, contrary to the hypothesis. Thus $H=M$. Any other
maximal subgroup containing the pair would contain $H=M$ and equal it.
\end{proof}

For $G=O'N$ and $(p,q)=(8,15)$ or $(15,16)$, the criterion applies
as follows.
The exact character evaluations give $x_M(p,q)=0$ for eleven of the
thirteen maximal conjugacy classes of $G$. The only remaining classes
are the two $L_2(31)$ classes, rows 15 and 16 of the roster. This zero
test uses the complete character calculation, not merely divisibility
of maximal-subgroup orders.

For each of the two embeddings, the complete maximal-subgroup list
inside $L_2(31)$ has orders
\[
465,\quad60,\quad60,\quad32,\quad30,\quad24,\quad24.
\]
None is divisible by $\operatorname{lcm}(p,q)$, which is 120 or 240.
Every proper pair therefore generates one of these two maximal classes.

Let $e_G(p,q)$ be the number of generating pairs. The proposition and
self-normalization of maximal subgroups in the simple group $G$ give
\[
e_G(p,q)=x_G(p,q)-2[G:L_2(31)]x_{L_2(31)}(p,q).
\]
The two acquired character inputs have equal values. Automorphisms
act freely on generating pairs, so division by
$|\operatorname{Aut}(G)|=921631011840$ gives the required counts:
\begin{center}\begin{tabular}{rrrrr}\toprule
$p$&$q$&$x_G/|\operatorname{Aut}G|$&$(x_G-e_G)/|\operatorname{Aut}G|$&$n_G$\\\midrule
8 & 15 & $3976$ & $16$ & 3960\\
15 & 16 & $15840$ & $32$ & 15808\\
\bottomrule\end{tabular}\end{center}

The intermediate columns are normalized numbers of pairs, and are
not asserted to count orbits of nongenerating pairs. There is no
additional division for duality. Both exact results agree with the
fixed original Table~1. They were not determined by the earlier
accepted bounds or the other exact-column arguments.

This is a new application of the shared subtraction argument, using
borrowed subgroup and character data; it is not a new classification.
It proves no further chirality assertion and changes no principal
refutation count. No new finite group acquisition is used.

\begin{thebibliography}{9}
\bibitem{CL} T. Connor and D. Leemans, \emph{Algorithmic enumeration
of regular maps}, Ars Math. Contemp. 10 (2016), 211--222.
\bibitem{roster} T. Connor and D. Leemans, \emph{The subgroup lattice
of $O'N$}, fixed 23-page roster, rows 15 and 16 and their maximal
subgroups. The exact edition and native pages are bound separately.
\end{thebibliography}
\end{document}
